Before we can talk about complexity and structure, we need a language to talk about uncertainty. That language is Information Theory.

\subsection{Entropy: A Measure of Surprise}
The foundational quantity is Shannon Entropy, denoted $H$. It measures the average uncertainty or "surprise" inherent in a random variable.

If we have a random variable $X$ that can take values $x$ from an alphabet $\mathcal{A}$ with probability $P(x)$, the entropy is:
\begin{equation}
    H[X] = - \sum_{x \in \mathcal{A}} P(x) \log_2 P(x)
\end{equation}
We measure entropy in \textbf{bits}.

\begin{sidebar}{Why Bits?}
    Shannon chose base 2 for his logarithms, giving us the \textit{bit} (binary digit). A bit represents a choice between two equally likely alternatives. However, other bases are possible:
    \begin{itemize}
        \item Base $e$ gives the \textit{nat}, used in thermodynamics \cite{shannon1948mathematical}.
        \item Base 10 gives the \textit{ban} or \textit{hartley} \cite{hartley1928transmission}.
    \end{itemize}
    We use bits because they map intuitively to digital logic.
\end{sidebar}

\begin{intuition}
    Think of entropy as the number of yes/no questions you need to ask on average to determine the value of $X$. If a coin is fair, you need 1 question ("Is it Heads?"). $H = 1$ bit. If the coin is rigged to always be Heads, you need 0 questions. $H = 0$ bits.
\end{intuition}

\subsection{Joint and Conditional Entropy}
What if we have two variables, $X$ and $Y$?
\begin{itemize}
    \item \textbf{Joint Entropy} $H[X, Y]$: Uncertainty about the pair $(X, Y)$.
    \item \textbf{Conditional Entropy} $H[X|Y]$: Uncertainty about $X$, given that we already know $Y$.
\end{itemize}
\begin{equation}
    H[X|Y] = H[X, Y] - H[Y]
\end{equation}
"The uncertainty about $X$ given $Y$ is the total uncertainty minus the uncertainty we resolved by learning $Y$."

\subsection{Mutual Information}
Perhaps the most important quantity for us is \textbf{Mutual Information} ($I$). It measures how much information $Y$ provides about $X$ (and vice versa).
\begin{equation}
    I[X; Y] = H[X] - H[X|Y]
\end{equation}
This is the reduction in uncertainty about $X$ gained by observing $Y$. If $X$ and $Y$ are independent, knowing $Y$ tells us nothing about $X$, so $I=0$.
